\scanpage{135}
\uline{de dégénérescence}, \uline{alors} (219) \uline{est un}\deleted{e}
\deleted{\unsure{équiv.}}\uline{isomorphisme}, de nature ess.\textsuperscript{t} triviale.

On est donc réduit à regarder le
cas d'une application \uline{injective} $\alpha$ (``opération
face''), et le \inserted{seul} cas qui \deleted{\unsure{pose}} pose problème
est bien sûr celui où $\alpha(0)\ne0$. Donc
\[
\tag{221}
\left\{\begin{aligned}
d_*&=\{d_0>d_1\mathinner{\cdot\,{-}\,{-}}>d_r\},\quad\text{et si}\\
d'_*&=\{d'_0>d'_1>\mathinner{\cdot\,{-}}>d'_{r'}\}\text{ est }
\text{\deleted{\unsure{une sous-suite}}}\inserted{\text{contenu dans }d_*},
\end{aligned}\right.
\]
\[
d'_0\ne d_0\text{ ie. }d'_0<d_0,\quad\text{il faut}
\]
définir l'hom
\[
(222)
\]
\diagram{\begin{tikzcd}[ampersand replacement=\&]
\Pi_1P^*_{d_0\deleted{\unsure{\geq}}d_1,\mathinner{{-}\,{-}},d_r}\arrow[r]\&
\Pi_1P^*_{d'_0,\mathinner{-},d'_r}
\end{tikzcd}}
\margin{attention, on a subrepticement chang\unsure{é}
de notations, en interver-tissant $d_*$ et $d'_*$ !}
\margin{on se réduit au cas
de $d_0,d_1,\mathinner{-},d_r$ strict.\textsuperscript{t}
\unsure{décroissant|décroissants} -- auquel on
parvient aussitôt. --
En fait, on est ramené au cas où $r'=r-1$, et où il
s'agit de l'opération ``face de codim $1$'' pour laquelle
$\alpha(0)\ne0$, donc correspondant :
\diagram{\begin{tikzcd}[ampersand replacement=\&]
D_{d_0,d_1,\mathinner{-},d_r}\arrow[r]\&D_{d_1,\mathinner{{-}\,{-}},d_r}
\end{tikzcd}}}

Notons que l'on a
\[
(223)
\]
\diagram{\begin{tikzcd}[ampersand replacement=\&]
D^{\deleted{*}}_{d_0,\mathinner{{-}\,{-}},d_r}\arrow[r]\&D_{d'_0,\mathinner{-},d'_{r'}}
\end{tikzcd}}
(mais \uline{non} un morphisme entre les
espaces avec $*$ !), qui est une immersion \inserted{locale}
d'espaces analytiques lisses, de codim
$d_0-d'_0$, \deleted{qui}\inserted{\unsure{ainsi|et}} s'insère dans le diagramme
\margin{134}
